% commentaire.tex — f175v, vv. 1082–1127 (conversion initiale des travaux antérieurs)
\comm{1082}{ὀτότυζ᾽ --- ὀτότυζ᾽.}{Signalons à l'attention du lecteur -- que j'espère exégète à ses heures -- qu'Aristophane invente le gentilé Ὀτοτύξιοι au vers 1043 des \italique{Oiseaux} par jeu de consonance avec Ὀλοφύξιοι (v. 1042), \italique{habitants d'Olophyxos}. Y a-t-il ici un jeu analogue ?}
\comm{1085}{σί κακόν;}{Commentaire du scholiaste :\citComm{ἀντὶ τοῦ `σοὶ κακόν'.}}
\comm{1089}{κακκασκι μοι;}{$=$ καὶ ἐγχάσκει μοι. Cf. \italique{Guêpes}, vers 721 et le scholiaste : καταγελᾷς μοι.}
\comm{1098–102}{ὦ θεοὶ ... κάρα κομίζων.}{La scholie à propos de la tirade de Persée\nc{Fr. 9 de l'\italique{Andromède} dans l'édition Jouan-van Looy $=$ fr. 124 Kannicht.} :\citComm{εἰς Περσέα ἐξ Ἀνδρομέδας τρία τὰ πρῶτα. καὶ λοιπὸν (i.e. vv. 1101-2) ἐπέζευξεν τὰ ἑξῆς.}}
\comm{1102–3}{τί λέγι; τὴ Γόργος πέρι / τὸ γραμματέο σὺ τὴ κεπαλή;}{$=$ τί λέγεις; τοῦ Γόργου φέρεις / τοῦ γραμματέως σὺ τὴν κεφαλήν;\\
Et la scholie (plus que douteuse) :\citComm{ὁ δὲ Γόργος γραμματεὺς, ἀλλὰ καὶ βάρβαρος.}Il n'y a en grec qu'une différence d'accent entre Γοργών (la Gorgone) et Γόργων (Gorgon).}
\comm{1105}{ἔα· τίν᾽ ὄχθον τόνδ᾽ ὁρῶ καὶ παρθένον.}{On connaît le texte original d'Euripide grâce à une scholie\nc{Dixième fragment de l'\italique{Andromède} dans l'édition Jouan-van Looy $=$ fr. 125, 1-2 Kannicht. La scholie rapporte les seuls deux premiers vers ; le fragment a pu être complété en le recoupant avec d'autres témoignages.} :\citComm{πάλιν ἐξ Ἀνδρομέδας·\\
ἔα, τίν᾽ ὄχθον τόνδ᾽ ὁρῶ περίρρυτον\\
ἀφρῷ θαλάσσης παρθένου τ᾽ εἰκώ τίνα\\
$[$ ἐξ αὐτομόρφων λαΐνων τυκισμάτων,
σοφῆς ἄγαλμα χειρός; $]$ }}
\comm{1106}{ναῦν ὅπως ὡρμισμένην;}{Quoique le scholiaste ne l'indique pas, le vers 1106 est probablement lui aussi repris de l'\italique{Andromède} d'Euripide. On trouve en tout cas la même image chez l'auteur d'\italique{Hercule Furieux} (vers 1094-1097), Hercule dont la situation était analogue à celle d'Andromède :\citComm{τί δεσμοῖς ναῦς ὅπως ὡρμισμένος\\
νεανίαν θώρακα καὶ βραχίονα,\\
πρὸς ἡμιθραύστῳ λαΐνῳ τυκίσματι\\
ἧμαι, νεκροῖσι γείτονας θάκους ἔχων;}}
